\begin{proof}[Proof of \cref{obj:lem:arbitrary-family-lower-obstruction}]
\begin{enumerate}
\item Fix the family, evaluation point, sample size, and positive constants in the first assertion. Define the constant estimator
\[
T_{\mathrm{known}}(\mathcal D_n):=\mu^\circ(x_0).
\]
It is measurable for every \(n\), including \(n=0\). For every \(P\in\mathfrak P\) and every sample,
\[
\left|T_{\mathrm{known}}(\mathcal D_n)-\mu_P(x_0)\right|
=\left|\mu^\circ(x_0)-\mu^\circ(x_0)\right|=0.
\]
Since \(c r_n>0\), its miss event is empty. Consequently, in the extended nonnegative reals,
\[
0\le
\inf_{T\ \mathrm{measurable}}\sup_{P\in\mathfrak P}
P^{\otimes n}\!\left\{
|T(\mathcal D_n)-\mu_P(x_0)|>c r_n
\right\}
\le
\sup_{P\in\mathfrak P}P^{\otimes n}(\varnothing)
=0.
\]
This proves the first assertion. % lean: CausalSmith.Stat.WeakOverlap.knownRegression_minimax_miss_zero

\item For the explicit witness, define the cube and component measures by
\[
\mathcal S:=[-1,1]^d,\qquad
\lambda_d:=\text{Lebesgue measure on }\mathbb R^d,
\]
\[
\Pi_X:=2^{-d}\lambda_d\!\restriction_{\mathcal S},
\qquad
\Pi_A:=\operatorname{Bernoulli}(1/2),
\qquad
\Pi_Y:=N(0,1),
\qquad
P^\circ:=\Pi_X\otimes(\Pi_A\otimes\Pi_Y).
\]
The cube-volume formula gives
\[
\lambda_d(\mathcal S)=\prod_{i=1}^d 2=2^d,
\qquad
\Pi_X(\mathbb R^d)=2^{-d}2^d=1.
\]
Also,
\[
\Pi_A(\{0,1\})=\frac12+\frac12=1,
\qquad
\Pi_Y(\mathbb R)=1.
\]
Thus \(P^\circ\) is a probability law with the independent coordinates specified in the statement. Its covariate marginal is
\[
P_X^\circ=\Pi_X.
\]
Select the treated regression, control regression, and propensity, respectively, as
\[
\mu_{P^\circ}(x):=0,\qquad
\mu_{0,P^\circ}(x):=0,\qquad
e_{P^\circ}(x):=\frac12
\quad (x\in\mathbb R^d).
\]
% lean: CausalSmith.Stat.WeakOverlap.dornCanonicalLaw_covariate

\item The Gaussian fourth absolute moment is finite. Define
\[
M_{\mathrm{mom}}
:=
\max\left\{
1,\left|\int_{\mathbb R}|y|^4\,\Pi_Y(dy)\right|
\right\}.
\]
Then
\[
1\le M_{\mathrm{mom}},
\qquad
\int_{\mathbb R}|y|^4\,\Pi_Y(dy)
\le
\left|\int_{\mathbb R}|y|^4\,\Pi_Y(dy)\right|
\le M_{\mathrm{mom}}.
\]
Monotonicity of powers in the exponent for a base at least one yields
\[
M_{\mathrm{mom}}
=M_{\mathrm{mom}}^1
\le M_{\mathrm{mom}}^4.
\]
Hence
\[
\int_{\mathbb R}|y|^4\,\Pi_Y(dy)
\le M_{\mathrm{mom}}^4.
\]
Choose the common constants
\[
q:=4,\qquad M:=M_{\mathrm{mom}},\qquad C:=2^{\gamma-1}.
\]
They satisfy \(q>3\), \(M>0\), and \(C>0\). % lean: CausalSmith.Stat.WeakOverlap.arbitraryFamily_lower_obstruction

\item We verify the remaining model restrictions with these constants. The product law gives
\[
\mathcal L_{P^\circ}(X,A)=\Pi_X\otimes\Pi_A.
\]
This is the joint law obtained from \(\Pi_X\) and the constant conditional distribution \(\Pi_A\). Uniqueness of conditional distributions therefore gives
\[
\mathcal L_{P^\circ}(A\mid X=x)=\Pi_A,
\qquad
P^\circ(A=1\mid X=x)=\frac12
\quad\text{for }P_X^\circ\text{-almost every }x.
\]
Thus the selected propensity is a measurable version, with
\[
0<e_{P^\circ}(x)=\frac12<1
\quad (x\in\mathbb R^d).
\]

Similarly, reassociating the product law gives
\[
\mathcal L_{P^\circ}((X,A),Y)
=(\Pi_X\otimes\Pi_A)\otimes\Pi_Y.
\]
Uniqueness of conditional distributions identifies the conditional outcome law with \(\Pi_Y\) for \((\Pi_X\otimes\Pi_A)\)-almost every \((x,a)\). Fubini's theorem gives this identity for \(\Pi_A\)-almost every \(a\), for \(\Pi_X\)-almost every \(x\). Since
\[
\Pi_A(\{0\})=\Pi_A(\{1\})=\frac12>0,
\]
both arms satisfy, simultaneously,
\[
\mathcal L_{P^\circ}(Y\mid X=x,A=a)=\Pi_Y
\quad(a\in\{0,1\}),
\qquad
P_X^\circ\text{-almost everywhere}.
\]
In particular, the treated conditional absolute fourth moment is integrable and
\[
\mathbb E_{P^\circ}\!\left[|Y|^4\mid X=x,A=1\right]
=
\int_{\mathbb R}|y|^4\,\Pi_Y(dy)
\le M^4
\quad\text{for }P_X^\circ\text{-almost every }x.
\]
The regression variance satisfies
\[
\int_{\mathbb R^d}
\left(
\mu_{P^\circ}(x)
-\int_{\mathbb R^d}\mu_{P^\circ}(y)\,P_X^\circ(dy)
\right)^2P_X^\circ(dx)
=0\le M.
\]

Both selected regressions have continuous derivatives of every order, and for every multi-index \(\alpha\) and every \(x,y\in\mathcal S\),
\[
D^\alpha\mu_{P^\circ}(x)
=D^\alpha\mu_{0,P^\circ}(x)=0,
\]
\[
|D^\alpha\mu_{P^\circ}(x)-D^\alpha\mu_{P^\circ}(y)|
=
|D^\alpha\mu_{0,P^\circ}(x)-D^\alpha\mu_{0,P^\circ}(y)|
=0.
\]
Their continuous boundary traces are also zero. Every defining top-derivative difference therefore vanishes, while its prescribed modulus bound is nonnegative. Thus
\[
\mu_{P^\circ},\mu_{0,P^\circ}\in\Sigma^{\mathrm D}(\beta,L_0).
\]

For the global overlap bound, fix \(t\in[0,1]\). If \(t\ge1/2\), monotonicity of the power with positive exponent \(\gamma-1\) gives
\[
(1/2)^{\gamma-1}\le t^{\gamma-1},
\qquad
2^{\gamma-1}(1/2)^{\gamma-1}=1.
\]
The propensity event is then the whole covariate space, so
\[
P_X^\circ\{e_{P^\circ}(X)\le t\}
=1
\le 2^{\gamma-1}t^{\gamma-1}
=Ct^{\gamma-1}.
\]
If \(t<1/2\), that event is empty, and
\[
P_X^\circ\{e_{P^\circ}(X)\le t\}
=0\le Ct^{\gamma-1}.
\]
Finally, the conditional outcome identities already established give Gaussian distributions in both arms with their selected zero means and variance one:
\[
\mathcal L_{P^\circ}(Y\mid X=x,A=1)
=N(\mu_{P^\circ}(x),1),
\qquad
\mathcal L_{P^\circ}(Y\mid X=x,A=0)
=N(\mu_{0,P^\circ}(x),1)
\]
for \(P_X^\circ\)-almost every \(x\). % lean: CausalSmith.Stat.WeakOverlap.arbitraryFamily_lower_obstruction

\item Define the selected singleton family and common local constants by
\[
\mathfrak P^\circ:=\{P^\circ\},
\qquad
\rho:=1,\qquad \nu:=\frac12,\qquad h_0:=1.
\]
The family is nonempty, its selected propensity has the required properties, and
\[
P_X^\circ(\mathcal S)=2^{-d}\lambda_d(\mathcal S)=1.
\]
We verify the local ratio in \cref{obj:def:dorn-a3}. For \(x_0\in\mathcal S\) and \(h>0\), define
\[
\mathcal B_h(x_0)
:=
\left\{x\in\mathbb R^d:
\sum_{i=1}^d(x_i-x_{0,i})^2\le h^2\right\},
\qquad
\mathcal O_h(x_0)
:=
\left\{x\in\mathbb R^d:
\sum_{i=1}^d(x_i-x_{0,i})^2<h^2\right\}.
\]
Since \(x_0\in\mathcal B_h(x_0)\cap\mathcal S\), this intersection is nonempty, and the constant selected propensity has supremum
\[
\sup_{x\in\mathcal B_h(x_0)\cap\mathcal S}e_{P^\circ}(x)
=\frac12.
\]

The denominator in the local ratio is positive, including when \(x_0\) lies on the cube boundary. Indeed, the coordinatewise closure identity
\[
\overline{(-1,1)}=[-1,1]
\]
and the finite-product closure and interior identities give
\[
\mathcal S\subseteq\overline{\operatorname{int}(\mathcal S)}.
\]
The set \(\mathcal O_h(x_0)\) is open and contains \(x_0\), because \(0<h^2\). It therefore meets \(\operatorname{int}(\mathcal S)\). Their intersection is a nonempty open set, so its Lebesgue measure is positive. Since
\[
\mathcal O_h(x_0)\cap\operatorname{int}(\mathcal S)
\subseteq
\mathcal B_h(x_0)\cap\mathcal S,
\]
we obtain
\[
0<
\lambda_d\!\left(
\mathcal O_h(x_0)\cap\operatorname{int}(\mathcal S)
\right)
\le
\lambda_d\!\left(\mathcal B_h(x_0)\cap\mathcal S\right).
\]
Consequently,
\[
0<
P_X^\circ(\mathcal B_h(x_0))
=
2^{-d}\lambda_d\!\left(\mathcal B_h(x_0)\cap\mathcal S\right)
\le1.
\]
Because the propensity is identically \(1/2\) on all of \(\mathbb R^d\),
\[
\left\{x:
e_{P^\circ}(x)\ge
\rho\sup_{y\in\mathcal B_h(x_0)\cap\mathcal S}e_{P^\circ}(y)
\right\}
\cap\mathcal B_h(x_0)
=
\mathcal B_h(x_0).
\]
Thus, for every \(x_0\in\mathcal S\) and \(0<h\le h_0\),
\[
\frac{
P_X^\circ\!\left(
\left\{x:
e_{P^\circ}(x)\ge
\rho\sup_{y\in\mathcal B_h(x_0)\cap\mathcal S}e_{P^\circ}(y)
\right\}
\cap\mathcal B_h(x_0)
\right)}
{P_X^\circ(\mathcal B_h(x_0))}
=1>\frac12=\nu.
\]
This proves
\[
\mathsf A_3^{\mathrm{Dorn}}
\bigl(\mathfrak P^\circ,(e_P)_{P\in\mathfrak P^\circ}\bigr).
\]
% lean: CausalSmith.Stat.WeakOverlap.dornCube_ball2_volume_pos

\item The singleton has the common known treated regression
\[
\mu_P(x)=0
\quad(P\in\mathfrak P^\circ,\ x\in\mathbb R^d).
\]
Applying the constant-estimator argument above with this regression and \(x_0=0\) gives, for every \(n\in\mathbb N\) and every \(c,r_n>0\),
\[
\inf_{T\ \mathrm{measurable}}
\sup_{P\in\mathfrak P^\circ}
P^{\otimes n}
\left\{
|T(\mathcal D_n)-\mu_P(0)|>c r_n
\right\}
=0.
\]
% lean: CausalSmith.Stat.WeakOverlap.arbitraryFamily_lower_obstruction
\end{enumerate}
\end{proof}
